% makale_orbit_bastirma_tr_v2_uzun.tex — revize makalenin Türkçe, kısaltılmamış, açıklayıcı sürümü
% Derleme: xelatex (iki kez)
\documentclass[12pt,a4paper]{extarticle}
\makeatletter\renewcommand\normalsize{\@setfontsize\normalsize{13}{16.5}}\makeatother
\usepackage{fontspec}
\usepackage{polyglossia}
\setdefaultlanguage{turkish}
\setotherlanguage{english}
\usepackage{amsmath,amssymb,bm}
\usepackage{graphicx}
\usepackage{booktabs}
\usepackage{tabularx}
\usepackage{array}
\usepackage[margin=2.4cm]{geometry}
\usepackage{xcolor}
\usepackage{hyperref}
\usepackage{caption}
\usepackage{float}
\usepackage{enumitem}
\hypersetup{colorlinks=true,linkcolor=blue!50!black,citecolor=blue!50!black,urlcolor=blue!50!black}
\captionsetup{font=small,labelfont=bf}
\setlength{\parskip}{0.5em}
\setlength{\parindent}{0pt}
\newcommand{\sekil}[3][0.95]{\begin{figure}[H]\centering\includegraphics[width=#1\textwidth]{#2}\caption{#3}\end{figure}}
\newcommand{\tahmin}{\textit{(tahmin)}}
\newcommand{\EslestirilmisParagraf}{\textit{(Eşleştirilmiş test sonuçları eklenecek.)}}


\title{Ters Yönde Dönen Kapalı Yörüngelerle Kuadrupol Kaçıklığının Simetrik Bileşenine Erişim\\[4pt]
\large Revize makalenin Türkçe, ayrıntılı ve açıklayıcı sürümü}
\author{I. Abedrabbo, S. Hacıömeroğlu\\ \small İstinye Üniversitesi, İstanbul}
\date{Çalışma taslağı}

\begin{document}
\sloppy
\maketitle

\begin{abstract}
\noindent Proton elektrik dipol momenti (EDM) için önerilen depolama halkasında, kuadrupollerin birkaç
mikronluk kaçıklığı geometrik (Berry) faz yoluyla gerçek bir EDM'yi taklit eder; $10~\mu$m rms kaçıklıkta
sahte sinyal, $10^{-29}~e\cdot$cm hedef hassasiyetini bir ile üç mertebe aşar. Her odaklayan--dağıtan
kuadrupol çiftinin kaçıklığını simetrik ve antisimetrik bileşenlere ayırınca, simetrik bileşenin düşük
harmoniklerde kapalı yörüngeyi antisimetrik bileşene göre dört ile yedi kat daha zayıf bozduğu görülür; bu
yüzden tek demetle yörünge düzeltmesi onu büyük ölçüde yerinde bırakır. Bu çalışmada, sahte EDM'nin donmuş
spin enerjisinde ölçüldüğü semplektik demet ve spin izlemesiyle, iki ters dönen demetin kapalı
yörüngelerinin birlikte fitlenmesinin simetrik bileşeni geri bulduğunu gösteriyoruz: antisimetrik bileşen
iki yörüngenin farkında, simetrik bileşen toplamında görünür. Statik BPM ofsetleri toplamda simetrik
bileşenle ayırt edilemez olduğundan yöntem, demet tabanlı hizalama (BBA) referansına göre çalışan bir
sürüklenme izleyicisidir. Yöntemin geçerlilik alanını, aynı kanala düşen hatalara karşı nicel olarak
inceliyoruz: BPM hareketi, elektrostatik saptırıcı alanı ve roll hataları, düzeltici türü, korelasyonlu
kaçıklık ve bir haftalık işletim senaryosu. Kuadrupol içi dipol düzelticiler, saptırıcı trimleri ve doğru
modellenmiş BPM montajıyla, iki demetli izleme kuadrupol polarite çevirmesine gerek kalmadan sahte EDM'yi
BBA'lar arasında hedefin çok altında tutar. Saptırıcı roll'unun birinci mertebe etkisinin yalnız ortalama
roll'a bağlı olduğunu ve Omarov ve ark.'nın öngördüğü ortak düşey trimle EDM kanalına dokunmadan
bastırıldığını da gösteriyoruz.
\end{abstract}

\tableofcontents
\newpage

%======================================================================
\section{Giriş}
%======================================================================
\subsection{Donmuş spin yöntemi ve sahte EDM}
Temel bir parçacığın kalıcı elektrik dipol momenti hem parite (P) hem zaman tersinmesi (T) simetrilerini
bozar; CPT değişmezliği varsayılırsa CP ihlali demektir. Standart Model protonun EDM'si için yaklaşık
$10^{-31}~e\cdot$cm öngörür; bunun üzerindeki her gözlem Standart Model ötesi fizik demektir. Önerilen proton
EDM deneyi~\cite{srEDM,omarov} donmuş spin yöntemini kullanır: protonlar ``sihirli'' momentumda
($p\approx0.7$~GeV/$c$) depolanır. Bu momentumda yatay düzlemde spin momentumu izler; yani spin ile
momentum arasındaki açı turdan tura değişmez. EDM varsa saptırıcıların radyal elektrik alanıyla etkileşir
ve spini düzlemin dışına döndürür. Bu dönme polarimetreyle ölçülür~\cite{polarimeter}; hızı
\begin{equation}
f \equiv \frac{dS_y}{dt},\qquad d=10^{-29}~e\cdot\text{cm}\ \longleftrightarrow\ f\approx1~\text{nrad/s}.
\end{equation}
Manyetik dipol momentinin bu dönmeyi taklit eden etkileşimlerine \emph{sahte EDM} denir. Bu çalışma,
manyetik kuadrupollerin kaçıklığından doğan sahte EDM'yi inceler.

\paragraph{Mekanizma.} Kuadrupolün merkezinde manyetik alan sıfırdır. Kuadrupol kaçıksa demet,
kaçıklıkla orantılı bir dipol alanı görür. Düşey kaçıklık $dy$ radyal bir alan üretir ve spini düşey
düzlemde döndürür; bu, EDM sinyaliyle aynı düzlemdir. Yatay kaçıklık $dx$ düşey bir alan üretir ve spini
yatay düzlemde döndürür. Bu dönmelerin her biri tek başına kapalı yörünge boyunca ortalamada sıfıra
kapanır. Ama dik eksenler etrafındaki dönmeler birbiriyle değişmez (komütatif değildir): önce bir eksen
etrafında sonra diğeri etrafında dönmek, ters sıradakiyle aynı sonucu vermez. Bu yüzden her turda küçük bir
net düzlem dışı dönme kalır. Bu geometrik (Berry) faz~\cite{berry1984,aharonov1987,farley2004,haciomeroglu2019}
sahte EDM'yi yatay ve düşey kaçıklıklarda \emph{bilineer} yapar:
\begin{equation}
f_\text{sahte}\propto dx\cdot dy\quad\Rightarrow\quad|f_\text{sahte}|\propto\sigma^2,
\label{eq:sigma2}
\end{equation}
burada $\sigma$ rms kaçıklıktır. $\sigma=10~\mu$m'de sahte sinyal, kaçıklık desenine bağlı olarak hedefin
$10$--$10^3$ katıdır (Bölüm~\ref{sec:benchmark}).

\subsection{Halka}
Ref.~\cite{omarov}'daki simetrik-hibrit halkayı kullanıyoruz: 24 FODO hücresi; her hücrede elektrostatik
saptırıcılar ve manyetik kuadrupoller. Odaklayan (QF) ve dağıtan (QD) kuadrupollerin gradyan büyüklüğü
aynıdır ($g=0.2$~T/m), işaretleri zıttır. Hücre şöyledir: QF $\to$ sürüklenme $\to$ saptırıcı $\to$
sürüklenme $\to$ QD $\to$ sürüklenme $\to$ saptırıcı $\to$ sürüklenme. Saptırıcılar 95.5~m yarıçaplı
silindirik kapasitörlerdir; alan $E_0=4.4$~MV/m. Şekil~\ref{fig:optik} hücreyi ve beta fonksiyonlarını,
Tablo~\ref{tab:halka} ana parametreleri gösterir.

\sekil[0.6]{fig_v2_optics.png}{Bir FODO hücresi: eleman yerleşimi (alt) ve beta fonksiyonları.
Saptırıcıların düşey odaklaması sıfırdır ($n=1$); yatay odaklamaları ölçülen yatay tune ile dahil
edilmiştir. Hücre 24 kez tekrarlanır.\label{fig:optik}}

\begin{table}[H]\centering\small
\caption{Halka ve demet parametreleri (Ref.~\cite{omarov} ile uyumlu).}\label{tab:halka}
\begin{tabular}{ll}\toprule
Parametre & Değer \\\midrule
FODO hücresi sayısı & 24 \\
Halka yarıçapı / çevre & 95.5 m / 800 m \\
Demet momentumu & 0.7007 GeV/$c$ \\
Göreli $\beta$, $\gamma$ & 0.598, 1.248 \\
Tur süresi & 4.46 $\mu$s \\
Betatron tune'ları ($Q_x,Q_y$) & (2.64, 2.3) \\
QF / QD'de $\beta$, yatay & 65 m / 35 m \\
QF / QD'de $\beta$, düşey & 42 m / 76 m \\
Dispersiyon ($D^\text{max}$) & 34 m \\
Kayma faktörü & $-0.25$ \\\bottomrule
\end{tabular}
\end{table}

\subsection{Simetrik ve antisimetrik bileşenler}
Kuadrupoller $n=0,\dots,47$ diye numaralanır; $k=0,\dots,23$ numaralı hücrede QF $n=2k$, QD $n=2k+1$'dir.
Bir düzlemde $n$ numaralı kuadrupolün kaçıklığı $v[n]$ olsun. Her hücre için
\begin{equation}
\begin{split}
v^s[2k]&=v^s[2k+1]=\tfrac12\left(v[2k]+v[2k+1]\right),\\
v^a[2k]&=-v^a[2k+1]=\tfrac12\left(v[2k]-v[2k+1]\right)
\end{split}
\label{eq:symanti}
\end{equation}
tanımlarız; her $n$ için $v[n]=v^s[n]+v^a[n]$. Simetrik bileşen $v^s$ hücredeki QF ve QD'yi aynı yöne,
antisimetrik bileşen $v^a$ zıt yönlere kaydırır. Denklem~\eqref{eq:symanti}'nin tanımladığı izdüşümler
$P_s$ ve $P_a=I-P_s$ birbirine diktir; her biri 24 boyutlu bir alt uzayı kapsar. Aynı ayrışım yatay ($v_x$)
ve düşey ($v_y$) kaçıklıklara ayrı ayrı uygulanır.

\subsection{Kaçıklığın düzeltilmesi ve bu çalışmanın katkısı}
Ref.~\cite{omarov} kaçıklığı iki adımda düzeltir. \emph{Demet tabanlı hizalama} (BBA), her kuadrupolün
gradyanını değiştirip demeti, yörüngenin bu değişime tepki vermediği konuma, yani manyetik merkeze yönlendirir;
ac uyarımla hızlandırılabilir~\cite{marti} ve paralel yapılabilir~\cite{huang}. Kalan artık, dört
konfigürasyonun birleştirilmesiyle sadeleştirilir: saat yönünde (CW) ve tersinde (CCW) dönen demetler, her
biri iki kuadrupol polaritesinde ($g\to-g$) [Ref.~\cite{omarov}, Denk.~(C2)]. Bu güçlüdür ama ek ölçüm
süresi ister ve örgünün iki polarite arasında yüksek simetride olmasını gerektirir.

Tepki matrisinin tekil değer ayrışımına (SVD) dayanan standart yörünge düzeltmesi~\cite{chung1993,ziemann,mirza,wegscheider}
sürekli ve demeti rahatsız etmeden çalışır. Bu çalışmada şunları gösteriyoruz:
\begin{enumerate}[nosep]
\item Tek demetle yörünge düzeltmesi simetrik bileşenin çoğunu yerinde bırakır, çünkü bu bileşen kapalı
yörüngeyi zayıf bozar (Bölüm~\ref{sec:mekanizma}).
\item İki ters dönen demetin birlikte fitlenmesi onu geri bulur (Bölüm~\ref{sec:ikidemet}).
\item Statik BPM ofsetleri iki demetli fitte simetrik bileşenle ayırt edilemez; bu yüzden yöntem bir BBA
referansına göre çalışan sürüklenme izleyicisidir.
\item Yöntemin geçerlilik alanı, aynı kanala düşen hatalara karşı nicel olarak belirlenir
(Bölüm~\ref{sec:gecerlilik}) ve bir haftalık işletim senaryosunda sınanır (Bölüm~\ref{sec:zaman}).
\end{enumerate}

%======================================================================
\section{Simülasyon yöntemi}
%======================================================================
\subsection{Demet ve spin izlemesi}
Bütün sonuçlar, tek bir parçacığı ve spinini Bölüm~1.2'deki örgüde Gauss--Legendre semplektik
integratörüyle~\cite{hairer} izleyen bir kodla elde edildi; spin Thomas--BMT denklemini~\cite{bmt1959}
izler. Saptırıcı alanı içeride $E_r=E_0R_0/R$, dışarıda sıfırdır (sert kenar). Kuadrupol kaçıklıkları,
roll'ları ve gradyan hataları, saptırıcı alan hataları ve roll'ları ve dipol düzelticiler halka
parametreleridir. EDM bağlaşımı $d_p=10^{-29}~e\cdot$cm'ye ayarlanınca kod $dS_y/dt=9.8\times10^{-10}$~rad/s
verir ve iki demet yönü için işaret zıttır; bu beklenen sonuçtur. Sahte EDM sonuçlarının hepsinde EDM
bağlaşımı kapalıdır.

\subsection{Tepki matrisleri}
Yörünge tepki matrisi $R$, kuadrupol kaçıklığından BPM okumasına giden doğrusal haritadır ($48\times48$;
her düzlem ve her demet yönü için ayrı); izlemeyle elde edildi. Kuadrupol tekmesi demet yönüyle işaret
değiştirdiği için $R_\text{CW}\approx-R_\text{CCW}$; ortalama sütun korelasyonu iki düzlemde de $-0.95$'tir.
Tam $-1$ olmamasının nedeni, demeti büken kuvvetin bir kısmının yönle işaret değiştirmeyen elektrik kuvveti
olmasıdır. Analitik ifade
\begin{equation}
R_{ij}=(KL)_j\frac{\sqrt{\beta_i\beta_j}}{2\sin\pi Q}\cos(|\Delta\psi_{ij}|-\pi Q)
\label{eq:R}
\end{equation}
düşeyde izlenen matrislerle birkaç yüzde içinde uyuşur, yatayda saptırıcı odaklaması yüzünden uyuşmaz; bu
yüzden her yerde izlenen matrisler kullanıldı.

\subsection{Sahte EDM'nin donmuş spin enerjisinde ölçülmesi}\label{sec:donmus}
\paragraph{Neden kapalı yörünge üzerinden fırlatma?} Kapalı yörüngenin dışından fırlatılan parçacık
betatron salınımı yapar; bunun $S_y$'ye etkisi sahte EDM'den mertebelerce büyüktür. Bu yüzden (i) 4 boyutlu
kapalı yörünge $(x,x',y,y')$, turdan tura konum varyansını en küçük yaparak bulunur; (ii) parçacık spini
boylamsal olarak bu yörüngenin üzerinden fırlatılır; (iii) $S_y(t)=a+bt+\sum_k[c_k\cos\omega_kt+d_k\sin\omega_kt]$
fitlenir ve yalnız seküler eğim $b$ alınır.

\paragraph{Enerji seçimi.} Kapalı yörünge fırlatma konumunu ve açısını belirler, enerjiyi belirlemez. Sahte
EDM ise enerjiye çok duyarlıdır: $\sigma=10~\mu$m civarında bağıl momentum değişimi $\delta=10^{-7}$, $f$'yi
yaklaşık $290\times$hedef değiştirir; yatay spin presesyon hızı ise $d\omega_h/d\delta=-5.24\times10^6$~rad/s
ile değişir ve bu eğim kaçıklıktan bağımsızdır. Kaçık bir örgüde yatay spinin donduğu ($\omega_h=0$) enerji,
tasarımdaki sihirli enerjiden kaçıklıkla orantılı bir miktar kayar (burada $\sim10^{-7}$). Yatay spini
dondurmanın birinci mertebe koşulu sihirli momentuma ulaşmak, ikinci mertebe koşulu ise hizalama hataları
ölçüsünde sihirli momentumun biraz yanına kilitlenmektir. Deney $\omega_h=0$'ı, ölçülen yatay spin fazı
üzerinden geri beslemeyle tutar~\cite{jedi_phaselock}; biz de bu koşulu referans alıyoruz.

\paragraph{Prosedür.} Her sahte EDM iki enerjide ölçülür: $(f_0,\omega_0)$ ve
$\delta_1=-\omega_0/(d\omega_h/d\delta)$'daki $(f_1,\omega_1)$; sonra $\omega_h=0$'a ara değerlenir:
$f^*=f_0+(f_1-f_0)(0-\omega_0)/(\omega_1-\omega_0)$. İki büyüklük de $|\delta|\le10^{-6}$'da doğrusaldır;
pratikte $|\omega_1|\lesssim3\times10^{-4}$~rad/s, yani ikinci nokta zaten donmuş enerjidedir. İkinci nokta
uzak düşerse ($|\omega_1|>10^{-3}$~rad/s) ara değerlenen enerjide üçüncü bir ölçüm yapılır. Donmuş enerji
etrafında simetrik bir momentum dağılımı aynı ortalama sinyali verir, çünkü $f$, $\delta$'ya doğrusaldır
(ölçülen: $f(\pm10^{-7})$ ortalaması ile $f(0)$ aynı; $\pm10^{-6}$'da fark \%1).

\paragraph{Sabit momentumla fırlatmanın sorunu.} Parçacığı fırlatma noktasında $|p|=p_\text{sihirli}$ ile
fırlatmak, sonucu fırlatma noktasına bağımlı yapar. Sert kenarlı saptırıcılarda alan sınırda aniden başlar ve
potansiyel basamağı yoktur; elektrik alanı saptırıcı içinde radyal harekete göre iş yapar. Bu yüzden bozulmuş
bir kapalı yörünge boyunca momentum hücreden hücreye $\sim10^{-7}$ değişir ve her fırlatma noktası farklı bir
enerji seçer. Aynı kaçıklık desenini 6, 12 ve 18 hücre kaydırmak (yani fırlatma noktasını değiştirmek), sabit
momentumlu sonucu 1088'den 2729, $-456$ ve 95 kat hedefe götürür; donmuş spin sonucu ise 82.04, 81.90, 81.89 ve
81.94'te kalır. Dört donmuş spin çözümü aynı yörüngedir: hücreler arası momentum farkları, tek bir kapalı
yörünge boyunca doğrudan ölçülenlerle dört basamak aynıdır. İki prosedür arasındaki fark fırlatma anındaki
kapalı yörünge büyüklüğüyle büyür: iyi düzeltilmiş yörüngelerde ihmal edilebilir, düzeltilmemişlerde değil
(Ek~\ref{ek:firlatma}).

\subsection{Doğrulamalar}\label{sec:benchmark}
\begin{itemize}[nosep]
\item Kuadrupoller hizalıyken sonuç sayısal tabanda: tarama ayarlarıyla $\approx0.008\times$hedef, tam ayarla
$<0.003$.
\item EDM açıkken beklenen hız.
\item Bütün saptırıcılarda düzgün roll $\phi$: birinci mertebe $|f|=G\gamma\cdot2\pi\phi/T_\text{dev}$
(Bölüm~\ref{sec:sapt}), $10^{-5}$ bağıl doğrulukla ve iki demet yönünde aynı işaretle.
\item Rastgele kuadrupol kaçıklığında $|f|\propto\sigma^p$, $p=2.0002$ ($\sigma=2.5$--$10~\mu$m, üç desen):
geometrik faz.
\end{itemize}
Uyarı: $\sigma^2$ testi mutlak değeri doğrulamaz. Yanlış enerjide yapılan ölçümün hatası da $\sigma^2$ ile
ölçeklenir, çünkü hem enerji kayması hem $df/d\delta$ kaçıklıkla orantılıdır. Eski (sabit momentumlu)
prosedürle de $p\approx2$ bulunmasının nedeni budur.

%======================================================================
\section{Sahte EDM mekanizması}\label{sec:mekanizma}
%======================================================================
\subsection{İki bileşenin yörüngede görünürlüğü}
Kaçık bir kuadrupol demete $\theta\propto g\,v$ tekmesi verir. QF ile QD'nin gradyanları zıt olduğundan,
antisimetrik desen bir hücrenin iki kuadrupolünde demeti \emph{aynı} yöne, simetrik desen \emph{zıt} yönlere
iter. Düzgün odaklama (smooth focusing) yaklaşımında simetrik desen saf bir $k=24$ hücre harmoniğidir ve
yörünge tepkisi $G_k\propto1/|Q^2-k^2|$ düşük harmoniklere göre yaklaşık 100 kat zayıf olurdu.

Ölçülen tepki çok farklıdır [Şekil~\ref{fig:modlar}(a)]: hücre harmoniği $m$ olan bir kaçıklık deseni için
simetrik/antisimetrik yörünge kazancı oranı $m\le3$'te $0.15$--$0.25$'tir ve $m=12$'de bire çıkar. Nedeni, bir
tekmenin yörüngeye etkisinin $\sqrt{\beta}$ ile ağırlıklanması ve $\beta$'nın QF ile QD'de çok farklı olmasıdır
(Tablo~\ref{tab:halka}). Simetrik bir hücre deseninin iki tekmesi bu yüzden birbirini götürmez. Kalan net
tekmenin antisimetrik desenin net tekmesine oranı
\begin{equation}
\frac{|\sqrt{\beta_F}-\sqrt{\beta_D}|}{\sqrt{\beta_F}+\sqrt{\beta_D}}\approx0.15
\end{equation}
olup iki düzlemde de ölçülen düşük harmonik oranıyla uyuşur. Yani simetrik bileşen yörüngeye kör değildir;
hem yörüngenin hem de aşağıda görüleceği gibi spinin en duyarlı olduğu düşük harmoniklerde antisimetrik
bileşenden dört ile yedi kat daha zayıftır.

Aynı yapı $R$'nin SVD'sinde de görülür [Şekil~\ref{fig:modlar}(b)]: tekil değerler sıralandığında en
küçüklerine karşılık gelen sağ tekil vektörler ağırlıklı olarak simetriktir. Eşik $10^{-2}$ olan bir kesik SVD,
tek demet matrisinin 48 doğrultusundan 17'sini (yatay) ve 24'ünü (düşey) atar; bunlar çoğunlukla simetrik
doğrultulardır.

\sekil{fig_v2_modes.png}{İki bileşenin yörüngede görünürlüğü. (a) Hücre harmoniği $m$ olan simetrik ve
antisimetrik desenin kapalı yörünge kazançlarının oranı (izlenen tepki matrisleri, CCW). Noktalı: $\sqrt\beta$
tahmini; kesikli: düzgün odaklama tahmini. (b) Tek demet tepki matrisinin normalize tekil değerleri, karşılık
gelen sağ tekil vektörün simetrik içeriğine göre renklendirilmiş.\label{fig:modlar}}

\subsection{Sahte EDM bilineer bir fonksiyonel}
Denklem~\eqref{eq:sigma2} genel olarak $f=\sum_{ij}W_{ij}v_{x,i}v_{y,j}=F(v_x,v_y)$ biçimindedir; $W$ örgünün
belirlediği bir çekirdektir. İki deseni de ayırınca
\begin{equation}
f=F(v_x^s,v_y^s)+F(v_x^s,v_y^a)+F(v_x^a,v_y^s)+F(v_x^a,v_y^a).
\end{equation}
Antisimetrik içerik yörüngede olduğu gibi sahte EDM'de de baskındır, ama daha küçük bir farkla: yörüngedeki
görünürlüğüne göre simetrik bileşen daha zararlıdır. Bu yüzden tek demetle yörünge düzeltmesi $\sigma$'yı
azaltır ama geometrik fazın bastırılmasını garanti etmez.
\IfFileExists{fig_v2_channels.png}{\sekil{fig_v2_channels.png}{Sahte EDM'nin dört kanalı, donmuş spin
enerjisinde ayrı ayrı ölçülmüş (CCW).}}{\textit{(Kanalların donmuş spin enerjisinde ölçüldüğü şekil, yeniden
ölçüm bitince eklenecek.)}}

%======================================================================
\section{Düzeltme yöntemleri}
%======================================================================
\subsection{Demet tabanlı hizalama}
BBA kuadrupolleri tek tek hizalar ve global tepki matrisini ters çevirmez; bu yüzden iki bileşene aynı
davranır. Yaklaşık \%1 rastgele gradyan hatası olduğunda tek bir kuadrupolün gradyanını değiştirmek bütün
halkanın optiğini biraz değiştirir (``nefes alma'', Ek~\ref{ek:nefes}) ve sıfır noktasını kapalı yörünge
bozulmasıyla orantılı bir miktar kaydırır; yinelemek bu kaymayı giderir. Simülasyonumuzda sekiz geçiş, rms
simetrik kaçıklığı yatayda 4.3'ten 0.9~$\mu$m'ye, düşeyde 3.3'ten 0.06~$\mu$m'ye indirdi. (Bu geçişlerin sahte
EDM değerleri eski, sabit momentumlu prosedürle ölçüldü; donmuş spinle tekrarı uzun bir koşu olarak planlandı.)

\subsection{Tek demetle yörünge düzeltmesi}\label{sec:tekdemet}
Standart düzeltme kaçıklığı $\hat v=R^+y$ ile kestirir; $R^+$ kesik SVD sözde-tersidir. Simetrik bileşen için
sınırlaması koşullanmadır: simetrik doğrultular en küçük tekil değerlere sahiptir ve kesmenin ilk attığı
doğrultulardır. Sonuç bu yüzden eşiğe çok bağlıdır. Eşik en büyük tekil değerin $10^{-2}$'si olduğunda
5~$\mu$m'lik simetrik sürüklenme yalnız 3.9~$\mu$m'ye iner; $3\times10^{-3}$ ile 0.05~$\mu$m'ye iner (14~nm BPM
gürültüsü).

\paragraph{Eşik nasıl seçilir?} Gerçek kaçıklık bilinmeden: \emph{uyuşmazlık ilkesi}, fit artığı
$\|S\hat v-y\|$'nin gürültü düzeyine indiği en büyük eşiği seçer. Bu kural her kestirici için ``gerçek
kaçıklığı bilerek seçilmiş en iyi eşiğin'' artığını en kötü 1.7 kat içinde yeniden üretir. Model yanlışsa
(örneğin BPM montajı yanlış varsayılmışsa) fit gürültü düzeyine hiç inemez; bu, model hatasını veriden
gösteren yararlı bir işarettir. Her yerde bu kural kullanıldı.

\paragraph{Fark yaklaşımı.} Kuadrupol gradyanlarını modüle edip yörünge farkını almak BPM ofsetlerini
sadeleştirir, ama fark matrisinin koşullanması kötüdür ve nefes alma etkisi baskındır (Ek~\ref{ek:nefes},
\ref{ek:fark}).

\subsection{İki demetle yörünge düzeltmesi}\label{sec:ikidemet}
İki ters dönen demet birlikte depolanır ve aynı BPM'lerle, paket zamanlamasıyla ayrılarak okunur. Kuadrupol
tekmesi yönle işaret değiştirdiğinden kapalı yörüngenin antisimetrik kısmı iki demette neredeyse zıt, simetrik
kısmı ise bir hücrenin iki kuadrupolünün ters sırayla geçilmesinden gelen bir faz kaymasına kadar neredeyse
aynıdır (Şekil~\ref{fig:cwccw}). İki tepki matrisinin toplamı ve farkı cinsinden (Frobenius normları, düşey):
antisimetrik bileşen fark kanalını ($\tfrac12\|(R_\text{CW}-R_\text{CCW})P_a\|=45$) sürer, toplamı yalnız zayıf
($6.8$); simetrik bileşen ikisinde de görünür ($6.9$ ve $6.5$). Toplam kanalında iki bileşen eşit görünürlüktedir;
tek demette ise simetrik bileşen yedi kat büyük bir antisimetrik tepkinin altında kalır. Ek~\ref{ek:oyuncak}
iki boyutlu bir örnek verir.

\sekil{fig_orbit_cwccw.png}{$10~\mu$m rms kaçıklık için iki demetin düşey kapalı yörüngeleri. (a) Antisimetrik
desen: yörüngeler neredeyse zıt, toplamları (siyah) küçük bir artık. (b) Simetrik desen: yörüngeler bir faz
kaymasına kadar neredeyse aynı; toplam kalır.\label{fig:cwccw}}

\paragraph{Algoritma.} Yöntem bir BBA referansına göre çalışır:
\begin{enumerate}[nosep]
\item BBA'dan sonra iki demetin okumalarını ($y^\text{ref}_\text{CW}$, $y^\text{ref}_\text{CCW}$) ve düzeltici
ayarlarını ($c^\text{ref}$) kaydet.
\item Her adımda iki demetin okumalarını $T_\text{ort}$ boyunca ortala ve değişimleri
$\Delta y_b=y_b-y^\text{ref}_b$ ($b=$CW, CCW) oluştur; düzelticilerin bilinen yörünge etkisini geri ekle.
\item Kaçıklık sürüklenmesini iki demetli en küçük kareler fitiyle kestir:
\begin{equation}
\Delta\hat v=S^+\begin{pmatrix}\Delta y_\text{CW}\\\Delta y_\text{CCW}\end{pmatrix},\qquad
S=\begin{pmatrix}R_\text{CW}-\rho I\\R_\text{CCW}-\rho I\end{pmatrix};
\label{eq:fit}
\end{equation}
kesme eşiği uyuşmazlık ilkesiyle seçilir; $\rho$ BPM montajını tanımlar (Bölüm~\ref{sec:bpm}; laboratuvara
sabit BPM için $\rho=0$).
\item $\Delta\hat v$'yi kuadrupol içine yerleştirilmiş dipol düzelticilerle uygula (Bölüm~\ref{sec:duzeltici}).
\item Kestirilen simetrik sürüklenmeyi izle; beklenen birikim eşiği aşınca BBA'yı tekrarla.
\end{enumerate}
Statik BPM ofsetleri iki demete aynı girer, yani tamamen toplam kanalındadır ve simetrik kaçıklıkla ayırt
edilemez. Referans olmadan 100~$\mu$m rms ofset, hedefin yaklaşık 200 katı sahte EDM'ye yol açar. 2. adımdaki
farklarda statik ofset sadeleşir; geriye ofsetin son BBA'dan beri kuadrupol merkezine göre \emph{değişimi}
kalır; buna $\delta$ diyoruz (Bölüm~\ref{sec:bpm}).

\paragraph{Performans.} Şekil~\ref{fig:bastirma}, $\sigma=10~\mu$m rastgele kaçıklık, 14~nm BPM gürültüsü ve
sabit $10^{-2}$ eşik için yöntemleri karşılaştırır. Sahte EDM, deneyin gözlemlediği
$C=\tfrac12|f_\text{CW}-f_\text{CCW}|$ büyüklüğüyle verilir (gerçek EDM demet yönüyle işaret değiştirir).
Düzeltilmemiş makine medyan $49\times$hedef; tek demet düzeltmesi $12\times$; iki demet düzeltmesi
$0.0015\times$, yani sayısal tabanda. Tek demet sonucu daha düşük bir eşikle tabana iner
(Bölüm~\ref{sec:tekdemet}, Tablo~\ref{tab:bpm}); yani sabit eşikte iki demetin üstünlüğünün büyük kısmı
koşullanmadan gelir. Asıl avantaj BPM'ler hareket ettiğinde ortaya çıkar (Bölüm~\ref{sec:bpm}). Ayrıca iki
demetli düzeltme, referans şemadaki kuadrupol polarite çevirmesini gerektirmez; bu, ölçüm başına ikinci bir
depolamayı gereksiz kılabilir.

\sekil[0.6]{fig_v2_suppression.png}{$\sigma=10~\mu$m'de sahte EDM (beş desen, çubuk = medyan): düzeltmesiz,
tek demet ve iki demet düzeltmesinden sonra (14~nm BPM gürültüsü, eşik $10^{-2}$, sabit BPM'ler). Noktalı:
sayısal taban.\label{fig:bastirma}}

%======================================================================
\section{Geçerlilik alanı}\label{sec:gecerlilik}
%======================================================================
İki demetli fit simetrik bileşeni toplam kanalında okur. Aynı kanala düşen her hata onunla yarışır. Bu bölüm
bu hataları inceler. Aksi belirtilmedikçe senaryo: BBA'dan sonra 5~$\mu$m rms saf simetrik sürüklenme
(yöntemin yakalaması gereken bileşen), 14~nm BPM gürültüsü, 10 desen; sahte EDM değerleri hedefin katı
cinsinden $C$'dir.

\subsection{BPM modeli}\label{sec:bpm}
BPM demeti kendi elektriksel merkezine göre ölçer. Demet $b$'nin BBA'dan sonraki okuma değişimini
\begin{equation}
\Delta r_b=R_b\Delta v-u+n_b,\qquad u=\rho\,\Delta v+\delta
\end{equation}
yazıyoruz: $\rho=0$ laboratuvara sabit BPM, $\rho=1$ kuadrupolüne rijit bağlı BPM, $\delta$ ise BPM
elektriksel merkezi ile kuadrupol manyetik merkezi arasındaki göreli sürüklenme (mekanik ve elektronik). Uzun
dalga boylu yer hareketi kuadrupolleri ve yakınlarındaki BPM'leri birlikte taşır, yani $\rho=1$'e
yaklaştırır~\cite{rossbach}. Rijit bağlı bir BPM'in okuması, demetin kuadrupol merkezine göre konumudur; bu,
sahte EDM'yi belirleyen büyüklüğün ta kendisidir. Bu durumda tek demet bile simetrik bileşeni, $-\Delta v$
terimi sayesinde görür.

Dört kestirici karşılaştırıldı: sabit BPM varsayan tek demet (E1), $(R-\rho I)$ kullanan tek demet (E2), sabit
BPM varsayan iki demet (E3), $(R-\rho I)$ kullanan iki demet (E4, Denk.~\eqref{eq:fit}).

\begin{table}[H]\centering\small
\caption{5~$\mu$m simetrik sürüklenme için kalan simetrik kaçıklık (200 desen, medyan) ve sahte EDM
(10 desen, medyan ve en büyük). Düzeltmesiz $C=6.3$ (10 desenin 9'u hedefin üstünde).}\label{tab:bpm}
\begin{tabular}{llccc}\toprule
$\rho$ & kestirici & kalan [$\mu$m] & $C$ medyan & $C$ en büyük\\\midrule
0 & E1 (tek demet) & 0.051 & 0.008 & 0.035\\
0 & E3 (iki demet) & 0.017 & 0.005 & 0.012\\
0.5 & E1 & 4.50 & 1.8 & 13\\
0.5 & E3 & 3.30 & 0.013 & 0.05\\
1 & E1 & 4.88 & 3.5 & 19\\
1 & E2 & 0.115 & 0.006 & 0.017\\
1 & E3 & 4.39 & 0.19 & 0.33\\
1 & E4 & 0.012 & 0.003 & 0.011\\\bottomrule
\end{tabular}
\end{table}

\paragraph{Sonuçlar.}
\begin{itemize}
\item Sabit BPM'lerde tek demet ve iki demet fitleri tabana ulaşır. İki demetin kalan artıktaki üstünlüğü üç
kattır; makinede modelde olmayan \%0.2 rms gradyan hatası varsa 1.2 kata iner.
\item BPM'ler kuadrupolleri izliyor ama kestirici sabit varsayıyorsa (E3), simetrik sürüklenmenin çoğu kalır
(4.4~$\mu$m); yine de $C$ hedefin altındadır (0.19, en büyük 0.33). Neden: kalan kısım betatron tune'undan uzak,
spinin az duyarlı olduğu harmoniklerdedir; fit $m=2$--3'teki gücün \%76--78'ini giderir. Marj yalnız üç ile
beş kattır. $\rho=0.5$'te aynı kestirici $C=0.013$ verir.
\item Montajı bilen kestiriciler (E2, E4) tabana ulaşır. E4, $\rho$'daki \%10'luk hataya tabandan çıkmadan
dayanır; E2 ise $\hat\rho=0.9$'da altı desenin birinde hedefi aşar.
\end{itemize}

\paragraph{Rijit bağlı BPM'ler.} Geriye kalan hata göreli sürüklenme $\delta$'dır. $\delta$ BPM'den BPM'e
bağımsızken [Şekil~\ref{fig:bpm}(b)], iki demetli sahte EDM $\sigma_\delta=1~\mu$m'ye kadar sayısal tabanda kalır;
tek demetli $\sigma_\delta^2$ ile büyür ve $1~\mu$m'de sekiz desenin yedisinde hedefi aşar. İki demetli fitin
kestirim hatası $1.1\,\sigma_\delta$ ile büyür, ama gücünün \%93'ü halkanın düzgün ötelenmesi ve eğilmesindedir
($m=0,1$). Bunlar $(R-I)$'nin neredeyse sıfır doğrultularıdır: halka bir bütün olarak öteleniyorsa kapalı yörünge
de onunla ötelenir. Aynı nedenle spin bu doğrultulara duyarsızdır; demet yine kuadrupol merkezlerinden geçer.
Rijit BPM'in ``göremediği'' şey ile spinin ``umursamadığı'' şey aynıdır.

BPM'i kuadrupolün içine ya da üstüne monte etmek lineer hızlandırıcılarda yaygındır (TTF2/European XFEL,
CLIC)~\cite{ttf2,pacman}. Işık kaynakları bunları senkrotron ışımasının girder'ı ısıtması yüzünden ayırır; bu
gerekçe proton halkasında geçerli değildir. Modern ışık kaynaklarında BPM elektroniği için 0.1~$\mu$m/saat
sürüklenme~\cite{sls2}
ve bir haftada $\le1~\mu$m~\cite{aps_mba} şartları konur; dolayısıyla BBA'lar arasında
$\sigma_\delta\lesssim1~\mu$m bugünkü teknolojiyle uyumludur.

\sekil{fig_v2_bpm.png}{BPM modeli. (a) 5~$\mu$m simetrik sürüklenmenin düzeltilmesinden sonra sahte EDM, sabit
($\rho=0$) ve rijit ($\rho=1$) BPM'ler için, Tablo~\ref{tab:bpm}'deki kestiricilerle (10 desen; çubuk =
medyan). (b) Göreli BPM--kuadrupol sürüklenmesi $\sigma_\delta$ altında rijit BPM'ler (8 desen). Kesikli: hedef;
noktalı: sayısal taban.\label{fig:bpm}}

\subsection{Elektrostatik saptırıcı hataları}\label{sec:sapt}
Bir saptırıcının bağıl alan hatası $\varepsilon$ demeti yatayda saptırır; demet ekseni etrafındaki roll $\phi$
ise düşey bir alan ekler: alan vektörü $\phi$ kadar dönünce düşey bileşen $E_y=E_0\sin\phi\approx E_0\phi$ olur.
Tek bir saptırıcıda $\varepsilon=10^{-4}$ ve $\phi=0.1$~mrad kapalı yörüngeyi 418 ve 525~$\mu$m bozar; hakemin
$\sim400~\mu$m tahmini doğrulanır. Elektrik kuvveti demet yönüyle işaret değiştirmediği için bu yörüngeler iki
demette aynıdır (korelasyon $+0.95$) ve tamamen toplam kanalına düşer; orada simetrik kuadrupol sürüklenmesiyle
yarışırlar.

\paragraph{Yörünge.} Yalnız kuadrupol sütunları içeren bir fit, saptırıcı sürüklenmesini simetrik kaçıklık
sanır: $\varepsilon$'da $10^{-7}$ ya da $\phi$'de 0.1~$\mu$rad başına yaklaşık 1~$\mu$m [Şekil~\ref{fig:sapt}(a)].
Düzlem başına 48 saptırıcı sütunu eklemek problemi tam rank tutar (96 bilinmeyen, 96 okuma), ama koşul sayısını
yatayda 73'ten $5.6\times10^3$'e, düşeyde 118'den $7.9\times10^3$'e çıkarır. Nedeni, birkaç saptırıcı ve
kuadrupol deseninin neredeyse aynı yörüngeyi üretmesidir (asal açı kosinüsleri 4--6 doğrultuda 0.99'un üstünde;
ikisi de tune yakınındaki harmonikleri sürer). Sızıntı böylece ortadan kalkar: simetrik hata kusursuz modelle
0.036~$\mu$m'de kalır, \%0.2 gradyan hatasıyla $\varepsilon=10^{-5}$'te 0.4~$\mu$m'ye çıkar; bedeli gürültü
büyütmesinin iki katına çıkmasıdır.

\paragraph{Spin, birinci mertebe.} Roll, düşey bir elektrik tekmesi verir: $\theta_E=(2\pi/48)\phi$. Kapalı
yörünge bunu kuadrupollerdeki manyetik tekmelerle dengeler: $\sum\theta_B=-\sum\theta_E$. Sihirli momentumda
elektrik alanı spini momentumla aynı açıda döndürür, manyetik alan ise $1+G\gamma$ kat fazla. Sonuç, tur başına
$G\gamma\sum\theta_B$ kadar net bir spin dönmesidir:
\begin{equation}
f=G\gamma\,\frac{2\pi}{48}\,\frac{\sum_i\phi_i}{T_\text{dev}}.
\label{eq:roll1}
\end{equation}
İzleme bunu $10^{-5}$ bağıl doğrulukla verir ($3\times10^{-7}$ rms rastgele roll'da $1.26\times10^8\times$hedef).
Bu terim iki demette aynı işaretlidir.

\paragraph{Birinci mertebe terim yalnız ortalama roll'a bağlı.} Denklem~\eqref{eq:roll1} yalnız $\sum_i\phi_i$'yi
içerir. Bunu doğrudan sınadık: aynı rastgele desenin ortalaması çıkarıldığında ($\sum\phi_i=0$) sonuç
$1.26\times10^8$'den sayısal tabana (0.003--0.007) iner; dört başka rastgele desende ve iki permütasyonda da
$|f|\le0.022$. Yani incelenen ideal donmuş spin koşullarında, tek tek roll'ların birinci mertebe katkısı yalnız
ortalamalarına bağlıdır; sıfır ortalamalı dağılımın ayrıntısından bağımsızdır. Fiziksel olarak: donmuş spinde spin
her yerde momentuma paraleldir ve küçük dönmeler birinci mertebede toplanır; yerleşimin etkisi ikinci mertebede
girer.

\paragraph{Çift ve tek kanal.} İki demetin ölçümlerinden iki kanal kurulur: çift kanal
$S_+=\tfrac12(f_\text{CW}+f_\text{CCW})$ ve tek kanal $S_-=\tfrac12(f_\text{CW}-f_\text{CCW})=\pm C$. Gerçek EDM
yönle işaret değiştirdiği için tek kanaldadır. Ortalama roll'un birinci mertebe etkisi çift kanaldadır.
Tablo~\ref{tab:kanal} bunu ölçümlerle gösterir.

\begin{table}[H]\centering\small
\footnotesize\caption{Kanal ayrışımı (hedefin katı; donmuş spin; iki yön). ``Ortak trim'' bütün saptırıcılara aynı roll
düzeltmesidir, yani ortalama roll'un çıkarılması.}\label{tab:kanal}
\begin{tabular}{lcccc}\toprule
Etki & CW & CCW & çift $S_+$ & tek $S_-$\\\midrule
Gerçek EDM ($10^{-29}~e\cdot$cm) & $+0.989$ & $-0.982$ & 0.004 & \textbf{0.986}\\
EDM + rastgele roll ($3\times10^{-7}$ rms) & $-1.26\times10^8$ & $-1.26\times10^8$ & $-1.26\times10^8$ & \textbf{1.02}\\
EDM + aynı roll + ortak trim & $+0.997$ & $-0.987$ & 0.005 & \textbf{0.992}\\
Düzgün radyal $B$ ($10^{-11}$~T, saptırıcılarda) & 0.023 & 0.011 & 0.017 & 0.006\\
0.3~$\mu$m kuadrupol kaçıklığı & $-0.29$ & $-0.03$ & $-0.16$ & 0.13\\
aynı + sıfır ortalamalı roll (bireysel trim yok) & $+1.31$ & $-1.79$ & $-0.24$ & \textbf{1.55}\\\bottomrule
\end{tabular}
\end{table}

Tablodan çıkan üç sonuç:
\begin{enumerate}
\item \textbf{Ortak trim EDM'yi bozmaz.} $10^8$ katlık çift sinyalin yanında tek kanal EDM'ye eşit kalır; ortak
trim çift kanalı 0.005'e indirirken tek kanal değişmez. Trim, roll'un ortalamasıyla aynı fiziksel büyüklüktür
(düzgün düşey elektrik alan) ve yeni bir kanal açmaz. Bu, incelenen koşullarda gerçek EDM'nin ortak trim geri
beslemesinden etkilenmediğini gösterir; bütün sistematikler altında tek kanalın korunacağını göstermez.
\item \textbf{Radyal manyetik alanın birinci mertebe etkisi yoktur.} Radyal $B$'nin düşey tekmesini kuadrupoller
dengeler; ikisi de manyetiktir ve spin dönmeleri birbirini götürür. Elektrik roll'da dengeleyen manyetik tekme
spini $G\gamma$ kat fazla döndürdüğü için net terim kalır. Bu sonuç düzgün ve statik bir radyal alan içindir;
uzayda değişen kaçak alanlar için ayrıca incelenmelidir.
\item \textbf{İkinci mertebe bir çapraz terim tek kanaldadır.} Kapalı yörünge sıfır değilken, sıfır ortalamalı
roll'lar düşey bir kapalı yörünge oluşturur; bu, kuadrupol kaynaklı yatay yörüngeyle çarpılıp geometrik faz
üretir. 0.3~$\mu$m kuadrupol kaçıklığında $C$, 0.13'ten 1.55'e çıkar. Bu terim ortak trimle temizlenmez; bireysel
saptırıcı roll trimi gerekir (aşağıda).
\end{enumerate}

\paragraph{Ortak trim geri beslemesi.} Çift kanal, polarimetrenin zaten ölçtüğü ve EDM'den $10^8$ kat büyük
olabilen bir sinyaldir. Bu sinyali ortak bir düşey elektrik trimiyle (üst ve alt elektrotlar arasında küçük
ortak bir gerilim farkı) sıfırlamak bizim önerimiz \emph{değildir}; temel tasarımda zaten vardır.
Ref.~\cite{omarov} (Bölüm~III~B ve Tablo~IV), iki demette aynı yönlü düşey dönmenin $10^{-6}$~rad/s'nin altında
tutulacağını belirtir. Bizim katkımız bu trimin (i) roll'un birinci mertebe etkisini bütünüyle kaldırdığını,
(ii) EDM kanalına dokunmadığını göstermek ve (iii) aşağıdaki kaçak formülüyle trime nicel bir gereksinim
koymaktır. Neden gerekli? Gerçek
makinede iki demet tam simetrik değildir; örneğin iki demet ayrı ayrı donmuş enerjiye kilitlenir ve enerjileri
farklı olabilir. Roll kaynaklı sinyalin enerjiye duyarlılığı $d\ln f/d\delta=-2.9$ (CCW) ve $-2.5$ (CW); yani
asimetri $\epsilon\approx(2.5$--$2.9)\Delta\delta$. Kalan tek kanal katkısı bir çarpımdır:
\begin{equation}
C_\text{kalan}\approx|S_+|\cdot\epsilon.
\end{equation}
Örnek: $\Delta\delta=10^{-7}$ \tahmin, $\epsilon\approx3\times10^{-7}$. Çift kanal $1.26\times10^8$'de bırakılırsa
$C\approx40$; $10^{-4}$~rad/s'ye ($10^5\times$hedef) kilitlenirse $C\approx0.03$. Böylece ``CW/CCW simetrisi
$10^{-8}$ olmalı'' gereksinimi, ``çift kanal yaklaşık $10^{-4}$~rad/s'de tutulmalı'' gereksinimine dönüşür.
Omarov'un $10^{-6}$~rad/s düzeyi bu gereksinimin 100 kat altındadır: $\Delta\delta=10^{-7}$ için kaçak
$C\sim3\times10^{-4}$ olur ve $\Delta\delta\approx10^{-5}$'e kadar $C\lesssim0.03$ kalır \tahmin. Yani iki eşik
farklı sorulara cevaptır: $10^{-4}$~rad/s bizim kaçak mekanizmamızın \emph{izin verdiği üst sınır},
$10^{-6}$~rad/s Omarov'un \emph{çalışma hedefi}dir; tutarlılık için ikincinin birincinin altında olması
yeterlidir. Ortalama düşey elektrik alan, EDM halkalarının bilinen ana sistematiği olan ortalama
radyal manyetik alanın elektrik karşılığıdır; deney onu zaten kontrol etmek zorundadır. Geri beslemenin dinamik
davranışı (gürültü, gecikme, çözünürlük) henüz simüle edilmedi.

\paragraph{Spin, ikinci mertebe ve bireysel trimler.} Alan hataları ile roll'lar birlikte yatay ve düşey kapalı
yörüngeler üretir; çarpımları kendi başına bir geometrik faz doğurur: kuadrupoller kusursuzken bile
$\varepsilon,\phi\sim3\times10^{-7}$'de $C\approx1$. Yalnız kuadrupolleri düzeltmek bunu gidermez (\%0.2 gradyan
hatasıyla $C=0.75$ medyan, sekiz desenin üçü hedefin üstünde). Çözüm saptırıcıları kaynağında, fitin saptırıcı
kısmını kullanarak trimlemektir: alan hatasını saptırıcı gerilimiyle, roll'u üst ve alt elektrotlar arasındaki
küçük bir gerilim farkıyla. Tek tek saptırıcı hataları kötü kestirilir (\%25--80'i kalır), ama kötü kestirilen
doğrultular yörüngeye zaten az etki eden doğrultulardır: saptırıcı kaynaklı yörünge \%97--99 azalır ve sahte EDM
$C=0.010$'a (\%0.2 gradyan hatasıyla 0.016) iner.
\EslestirilmisParagraf

\paragraph{Saptırıcı kararlılığı.} Şekil~\ref{fig:sapt}(b) saptırıcı sürüklenmesinin genliğini tarar. Model
kusursuzken trimlenmiş sonuç $10^{-5}$'e kadar tabanda kalır, çünkü kestirim hatası yalnız BPM gürültüsüyle
belirlenir. \%0.2 gradyan hatasıyla genliğin karesiyle büyür ve $3$--$5\times10^{-6}$'da hedefe ulaşır. Yani trim
güncellemeleri arasında izin verilen saptırıcı sürüklenmesini tepki modelinin doğruluğu belirler.
Karşılaştırma için: bugünkü yüksek gerilim kaynakları $10^{-5}$ dalgalanma~\cite{cpedm_feas} ve $10^{-4}$ uzun
süreli kararlılık~\cite{fug} belirtir. Ortak mod değişimleri donmuş spin geri beslemesince emilir; yalnız trim
güncellemeleri arasındaki bağımsız yavaş sürüklenmeler önemlidir.

\sekil{fig_v2_deflector.png}{Saptırıcı hataları. (a) Saptırıcı alan sürüklenmesine karşı yeniden kurulan simetrik
kuadrupol sürüklenmesinin hatası (5~$\mu$m, yatay, 200 desen), yalnız kuadrupol sütunlu ve saptırıcı sütunları
eklenmiş fitler için. (b) Kuadrupol düzeltmesi ve saptırıcı triminden sonra sahte EDM, saptırıcı sürüklenme
genliğine karşı (4--8 desen).\label{fig:sapt}}

\subsection{Düzeltici türü}\label{sec:duzeltici}
Bölüm~\ref{sec:ikidemet}'teki yöntem, kaçıklığı oluştuğu yerde iptal eden bir düzeltici ister. Kuadrupolün
içine sarılmış bir dipol sargısı düzgün bir $\Delta B$ alanı ekler ve manyetik merkezi $\Delta B/g$ kadar
kaydırır: $g(y-d)+\Delta B=g(y-[d-\Delta B/g])$. $\Delta B=g\,\Delta\hat v$ seçilince kuadrupol iki demet ve spin
için hizalıymış gibi davranır; yukarıda varsayılan düzeltme budur (kuadrupol konumlandırıcıları eşdeğerdir).
Saptırıcılardaki düzelticiler bunu yapamaz:
\begin{itemize}
\item Oradaki manyetik düzelticiler iki demeti zıt yönlere iter ve yalnız fark kanalına etki eder; simetrik
sürüklenme ise iki demetin aynı yöne kaydırılmasını ister. BPM okumalarını sıfırlamak demeti kuadrupol
merkezlerinden geçirmez.
\item Elektrostatik düzelticiler iki demete aynı etki eder ve onları kuadrupol merkezlerinden geçirebilir
(0.1~$\mu$m içinde), ama gereken elektrik tekmeleri ($\varepsilon$ ve $\phi$'de $5$--$8\times10^{-6}$) yukarıdaki
birinci ve ikinci mertebe saptırıcı etkilerini üretir: okumaları sıfırlamada $C=7.6$, merkezlere yönlendirmede
$C=1.3$; ayrıca CW/CCW simetrisine tamamen dayanan $10^6$--$10^8\times$hedeflik aynı işaretli sinyaller.
\end{itemize}
Yalnız kuadrupol içi düzelticiler $C$'yi tabanda tutar.

\subsection{Korelasyonlu kaçıklık}
Ölçme hataları ve yer hareketi halka boyunca korelasyonlu kaçıklıklar üretir; iki kuadrupolün göreli yer
değiştirmesi aralarındaki mesafenin kareköküyle büyür. Bunu periyodik bir rastgele yürüyüşle modelledik:
$\langle(v_i-v_j)^2\rangle\propto|s_i-s_j|$, yani ATL yasasının ($\langle\Delta Y^2\rangle=ATL$~\cite{shiltsev})
uzaysal yapısı; zararsız rijit öteleme ve eğim çıkarıldı. Bu model belirli bir korelasyon senaryosudur; ortak
destek yapıları ve yerel mekanik hareketler gerçek makinede ek katkılar verebilir. Komşu kuadrupoller birlikte
hareket ettiği için simetrik bileşen baskınlaşır: simetrik/antisimetrik rms oranı 2.3 (varyansın \%84'ü),
korelasyonsuz hatalarda yaklaşık 1 (2000 desen). Yani korelasyon simetrik bileşeni daha önemli yapar; bu, onu
okuyan bir yöntemin gerekçesini güçlendirir. İki demetli düzeltme doğrusaldır ve desen istatistiğinden
etkilenmez: düzeltmeden sonra artık korelasyonsuz, rastgele yürüyüş ve karma modellerde aynı gürültü
düzeyindedir. Yer hareketi BPM'leri de taşıyorsa durum Bölüm~\ref{sec:bpm}'deki $\rho\approx1$ durumudur: sabit
BPM varsayan kestirici simetrik sürüklenmenin yalnız üçte birini geri bulur; $\rho=1$ kullanan kestirici
0.01~$\mu$m'ye kadar bulur.

%======================================================================
\section{İşletim senaryosu}\label{sec:zaman}
%======================================================================
Makineyi bir hafta boyunca 1 saatlik adımlarla izledik. Kuadrupoller ATL yasasıyla sürüklenir:
$A=10^{-5}~\mu$m$^2$/(s\,m), LEP'te ölçülen değere yakın~\cite{shiltsev}. Her BBA'dan sonra kalan kaçıklık
0.3~$\mu$m rms; BPM--kuadrupol göreli sürüklenmesi $\delta$, $0.2~\mu$m$/\sqrt{\text{gün}}$ ile büyür; beklenen
değeri 0.3~$\mu$m'ye ulaşınca BBA tekrarlanır (her $\approx54$~saatte). BPM okumalarında tur başına 10~$\mu$m
çözünürlük, 60~s ortalama ve 10~nm taban var. Karşılaştırılan durumlar: rijit BPM ve iki demetli (E4) ya da tek
demetli (E2) düzeltme; zeminle hareket eden ama sabit varsayılan BPM'ler (E3); yalnız periyodik BBA; düzeltme yok.

Sonuç (Şekil~\ref{fig:zaman}): iki demetli izleme sahte EDM'yi hafta boyunca BBA artığının belirlediği düzeyde
tutar (medyan $C=0.05$--$0.09$, 15 değerin en büyüğü 0.18; BBA'dan hemen sonra $C=0.06$). Tek demetli izleme ve
sabit BPM varsayımı hedefin çevresinde dolaşır (medyan en fazla 1.2 ve 1.3; 96. saatte beş desenin üçü ve
dördü hedefin üstünde). Yalnız periyodik BBA, iki kalibrasyon arasında $C\approx45$'e çıkar. İzleyici,
$3k\sigma_n$ büyüklüğünde bir simetrik sürüklenmeyi algılar; $k=0.75$ iki demetli simetrik kestirimin gürültü
kazancı, $\sigma_n$ okuma başına BPM gürültüsüdür; 10~nm tabanla bu $\approx0.02~\mu$m'dir. Sahte EDM $\sigma^2$
ile ölçeklendiği için, BBA artığı yaklaşık 0.4--0.5~$\mu$m'ye kadar kabul edilebilir \tahmin.

Sınırlamalar: spin sonuçları 5 desene dayanır; saptırıcı sürüklenmesi ve trimler bu senaryoya henüz dahil
değildir; BBA'nın kendi süresi ve gerçek zaman kazancı burada hesaplanmadı. Bunlar planlanan birleşik
senaryonun konusudur.

\sekil{fig_v2_time.png}{Bir haftalık işletim senaryosu. (a) Kalan simetrik kaçıklık (doğrusal model, 20 desen,
medyan ve \%16--84 bandı). (b) İzlemeyle sahte EDM (5 desen; çizgi = medyan). BBA 54, 108 ve 162. saatlerde
tekrarlanır.\label{fig:zaman}}

%======================================================================
\section{Sonuçlar}
%======================================================================
Kuadrupol kaçıklığından gelen sahte EDM, yatay ve düşey kaçıklık desenlerinin bilineer bir fonksiyonelidir.
Simetrik ve antisimetrik bileşenlere ayrıldığında simetrik bileşen, düşük harmoniklerde kapalı yörüngede dört
ile yedi kat daha az görünür (düzgün odaklama bastırmasından değil, QF ile QD arasındaki $\beta$ farkından);
sahte EDM'deki ağırlığı ise daha küçük bir oranda azalır. Tek demetle yörünge düzeltmesi, eşik dikkatle
seçilmedikçe onu büyük ölçüde yerinde bırakır.

İki ters dönen demet bileşenleri ayırır: antisimetrik olanı iki yörüngenin farkında, simetrik olanı toplamında.
Statik BPM ofsetleri toplamda simetrik bileşenle ayırt edilemediği için yöntem BBA kalibrasyonları arasında bir
sürüklenme izleyicisidir ve referans şemadaki kuadrupol polarite çevirmesine ihtiyaç duymaz. Geçerlilik alanını
aynı kanala düşen hatalar belirler:
\begin{itemize}
\item BPM modeli montajla uyuşmalıdır; rijit bağlı BPM'ler ve bunu bilen bir model, yöntemi en az 1~$\mu$m'lik
BPM--kuadrupol sürüklenmelerine ve montaj katsayısındaki \%10'luk hatalara karşı sağlam yapar;
\item saptırıcı sürüklenmeleri fite dahil edilmeli ve kaynağında trimlenmelidir; trimler arası izin verilen
sürüklenme \%0.2 gradyan hatasıyla $\sim3\times10^{-6}$'dır ve tepki modelinin doğruluğuyla ölçeklenir;
\item saptırıcı roll'unun birinci mertebe etkisi yalnız ortalama roll'a bağlıdır, çift kanaldadır ve
Ref.~\cite{omarov}'un ortak düşey trimiyle EDM kanalına dokunmadan bastırılır; roll ile kapalı yörüngenin ikinci mertebe çapraz
terimi tek kanaldadır ve bireysel saptırıcı trimi gerektirir;
\item düzeltmeler kuadrupol içi dipol düzelticilerle (ya da konumlandırıcılarla) uygulanmalıdır.
\end{itemize}
Bu koşullarda iki demetli izleme, simüle edilen bir haftalık işletim boyunca sahte EDM'yi BBA kalibrasyonları
arasında hedefin çok altında tutar. Son olarak, kaçık bir örgüde sahte EDM donmuş spin enerjisinde
hesaplanmalıdır; sabit momentumla yapılan hesap fırlatma noktasına bağlıdır ve mertebelerce farklı sonuç
verebilir.

%======================================================================
\appendix
\section{Donmuş spin hesabı ve fırlatma noktasına bağımlılık}\label{ek:firlatma}
Burada kullanılan desenler için sabit momentumla ($|p|=p_\text{sihirli}$, fırlatma noktasında) ve donmuş spin
enerjisinde hesaplanan sahte EDM'ler arasındaki fark, fırlatma anındaki kapalı yörünge büyüklüğüyle orantılıdır.
Şekil~\ref{fig:bastirma}'deki desenler için fark, iki demet düzeltmesinden sonra $\le0.02\times$hedef (rms
yörünge çarpımı $\approx0.01~\mu$m$^2$), tek demet düzeltmesinde düzeltilen demet için $\le0.2$
($\approx1~\mu$m$^2$), düzeltilmemiş demet ve makine için 1--3000'dir ($10^2$--$10^4~\mu$m$^2$). Hata iki demet
yönünde farklı olduğu için $C$'de sadeleşmez. Donmuş spin sonucu ara değerlemenin başlangıç kuralından
bağımsızdır: sabit momentum kuralından başlanınca dört fırlatma noktasında bulunan donmuş enerji, donmuş spin
yörüngesinin momentum profiliyle dört basamak aynıdır ve orada yapılan ölçüm 82.04 ve 81.87 verir; varsayılan
kuralla 82.04 ve 81.90.

\section{Optik nefes alma}\label{ek:nefes}
Bir kuadrupolün gradyanını değiştirmek bütün halkanın betatron fonksiyonlarını ve dolayısıyla her tepki matrisi
elemanını [Denk.~\eqref{eq:R}] biraz değiştirir. Bütün kuadrupoller farklı frekanslarda modüle edildiğinde bu
değişimler yeniden kurmaya hâkim olur.
\sekil[0.6]{fig_orbit_breathing.png}{Kuadrupoller modüle edilirken kaçıklıkların yeniden kurulması. Nefes alma
yapay olarak kaldırılırsa kurma çalışır; varken sinyal gömülür.}

\section{Yörünge farkının tersinmesi}\label{ek:fark}
İki gradyan ayarı arasındaki yörünge farkı statik ofsetleri sadeleştirir. Gradyan $\epsilon_g$ kadar değişirse
fark matrisi $\Delta R\approx\epsilon_g\,\partial R/\partial g$ ve $\|\Delta R^{-1}\|\approx\|R^{-1}\|/\epsilon_g$
olur; bu örgüde koşul sayısı $3.7\times10^4$'tür. En küçük tekil değerler simetrik doğrultulara ait olduğundan
tersinme seçici olarak simetrik bileşende başarısız olur; düzenlileştirme aynı doğrultuları atar.
\sekil[0.6]{fig_orbit_lockin.png}{Fark yaklaşımının $\beta$-beat'e karşı yeniden kurma hatası. Simetrik bileşenin
hatası antisimetriğinkinin yaklaşık 15 katı olur.}

\section{Ters dönen toplamın iki boyutlu gösterimi}\label{ek:oyuncak}
Bir düzlemdeki kaçıklığı $v=(a,s)^\top$ ile gösterelim: bir antisimetrik ve bir simetrik koordinat. Antisimetrik
olan büyük, simetrik olan küçük bir yörünge sürer: $R_\text{CW}=\mathrm{diag}(G,g)$, $G\gg g$. Demet yönünü
çevirmek antisimetrik yörüngeyi çevirir ama simetrik olanı (bir faz kaymasına kadar) çevirmez:
$R_\text{CCW}=\mathrm{diag}(-G,g)$. Her okuma, BPM'inin iki demete ortak statik ofsetini $b=(b_a,b_s)^\top$ taşır:
\begin{equation}
y_\text{CW}-y_\text{CCW}=(2Ga,\,0)^\top,\qquad y_\text{CW}+y_\text{CCW}=(2b_a,\,2gs+2b_s)^\top.
\end{equation}
Fark $a$'yı ofsetsiz ve iyi koşullu verir; toplam $s$'yi yalnız ofset $b_s$ ile birlikte ve küçük kazanç $g$
üzerinden verir. Statik simetrik kaçıklık ile statik ofset tek ölçümde ayırt edilemez; ama ofset toplamın
\emph{değişiminde} sadeleşir: $\Delta s=\Delta(y_\text{CW}+y_\text{CCW})/2g$. Yöntemin bir sürüklenme izleyicisi
olmasının nedeni budur. Halkada $v$ düzlem başına 48 elemanlıdır, her demet 48 okuma verir ve bileşenler
Denk.~\eqref{eq:fit}'in en küçük kareler çözümüyle bulunur. Gerçek halkada parite yaklaşıktır (demetler hücreyi
ters sırayla geçer); bu yüzden antisimetrik yörünge toplamda küçük bir artık bırakır [Şekil~\ref{fig:cwccw}(a)].

\begin{thebibliography}{99}\small
\bibitem{srEDM} V.~Anastassopoulos \emph{et al.}, Rev.\ Sci.\ Instrum.\ \textbf{87}, 115116 (2016).
\bibitem{omarov} Z.~Omarov \emph{et al.}, Phys.\ Rev.\ D \textbf{105}, 032001 (2022).
\bibitem{polarimeter} F.~Müller \emph{et al.} (JEDI), J.\ Instrum.\ \textbf{15}, P12005 (2020).
\bibitem{berry1984} M.~V.~Berry, Proc.\ R.\ Soc.\ Lond.\ A \textbf{392}, 45 (1984).
\bibitem{aharonov1987} Y.~Aharonov and J.~Anandan, Phys.\ Rev.\ Lett.\ \textbf{58}, 1593 (1987).
\bibitem{farley2004} F.~J.~M.~Farley \emph{et al.}, Phys.\ Rev.\ Lett.\ \textbf{93}, 052001 (2004).
\bibitem{haciomeroglu2019} S.~Hacıömeroğlu, Y.~F.~Orlov, Y.~K.~Semertzidis, Nucl.\ Instrum.\ Methods A
\textbf{921}, 12 (2019).
\bibitem{marti} Z.~Martí \emph{et al.}, Phys.\ Rev.\ Accel.\ Beams \textbf{23}, 012802 (2020).
\bibitem{huang} X.~Huang, arXiv:2203.14869.
\bibitem{chung1993} Y.~Chung, G.~Decker, K.~Evans, PAC 1993, p.~2263.
\bibitem{ziemann} V.~Ziemann and M.~Ziemann, arXiv:2104.05300.
\bibitem{mirza} S.~H.~Mirza \emph{et al.}, Phys.\ Rev.\ Accel.\ Beams \textbf{22}, 072804 (2019).
\bibitem{wegscheider} S.~Wegscheider \emph{et al.}, Phys.\ Rev.\ Accel.\ Beams \textbf{26}, 032803 (2023).
\bibitem{hairer} E.~Hairer, C.~Lubich, G.~Wanner, \emph{Geometric Numerical Integration} (Springer, 2006).
\bibitem{bmt1959} V.~Bargmann, L.~Michel, V.~L.~Telegdi, Phys.\ Rev.\ Lett.\ \textbf{2}, 435 (1959).
\bibitem{jedi_phaselock} N.~Hempelmann \emph{et al.} (JEDI), Phys.\ Rev.\ Lett.\ \textbf{119}, 014801 (2017).
\bibitem{rossbach} J.~Rossbach, Part.\ Accel.\ \textbf{23}, 121 (1989).
\bibitem{shiltsev} V.~Shiltsev, Phys.\ Rev.\ Lett.\ \textbf{104}, 238501 (2010).
\bibitem{ttf2} TTF2 BPM--kuadrupol gerili tel kalibrasyonu, LINAC 2004, TUP69.
\bibitem{pacman} CLIC/PACMAN BPM--kuadrupol ön hizalama, doktora tezi, Pisa Üniv.\ (2016).
\bibitem{sls2} ``Development of the SLS 2.0 BPM system'', IBIC 2023, MO3C03.
\bibitem{aps_mba} R.~Lill, N.~Sereno, B.~Yang, ``BPM stability studies for the APS MBA upgrade''.
\bibitem{cpedm_feas} F.~Abusaif \emph{et al.} (CPEDM), CERN-2021-003; arXiv:1912.07881.
\bibitem{fug} P.~Schmidt-Wellenburg, T.~Schmidt, FuG HCB 200~kV sunumu (2020).
\end{thebibliography}
\end{document}
